\documentclass{article}
\usepackage[utf8]{inputenc}
\usepackage[russian]{babel}
\usepackage{graphicx}
\usepackage{minted}
\usepackage{xcolor}
\usepackage[T2A]{fontenc}

\setminted{
    encoding=UTF-8,
    breaklines=true,
    frame=single,
    fontsize=\small,
}

\begin{document}
\begin{titlepage}
    \centering
    \vspace*{2cm}
    
    {\Large МИНИСТЕРСТВО НАУКИ И ВЫСШЕГО ОБРАЗОВАНИЯ РОССИЙСКОЙ ФЕДЕРАЦИИ
Федеральное государственное автономное образовательное учреждение высшего образования
“Санкт-Петербургский национальный исследовательский университет информационных технологий механики и оптики”
    \par}
    {\Large Мегафакультет: трансляционных информационных технологий
    \par}
    {\LargeФакультет: информационных технологий и программирования
    \par}
    \vspace{0.5cm}
    {\large Университет ИТМО\par}
    
    \vspace{1.5cm}
    {\Large\bfseries По дисциплине: “Инструментальные средства разработки ПО”
    \par}
    {\Large\bfseries Тема "Документирование"
    \par}
    
    \vspace{0.5cm}
    % {\Large Документирование\par}
    
    \vfill
    {\large Выполнил студент группы №M3121:\par}
    {\large Галочкин Тимофей Антонович\par}
    \vspace{0.5cm}
    {\large Работу проверила:\par}
    {\large Алексеева Милена Константиновна\par}
    \vspace{1.05cm}
    {\Large САНКТ-ПЕТЕРБУРГ\par}
    {\large 2026\par}
\end{titlepage}

\section{Цель работы}
    {\large - Научится писать, как отдельную документацию, так и внедрять её в исходный код\par}

\section{Документация к файлу mysterious\_file.py}
    % {\large \centering \boldsymbol{Уровень 1: Новичок (Базовый) — «Мой первый комментарий»}\par}
    % \vspace{0.5cm}
    

    \subsection{Функция bought\_products}
    
    Функция вычитает сумму покупки из денег пользователя.
    
    \begin{minted}{python}
    def bought_products(money, price_per_item, items_count):
        total = price_per_item * items_count
        if total > money:
            return "Недостаточно денег"
        return money - total
    \end{minted}
    
    \subsection*{Параметры}
        \begin{enumerate}
            \item \texttt{money (int)} - число денег у пользователя
            \item \texttt{price\_per\_item (float)} - цена одного товара
            \item \texttt{items\_count} (int) - число одинаковых товаров
        \end{enumerate}
        
        \subsection*{Возвращает}
        \begin{enumerate}
            \item \texttt{int} или \texttt{float} - оставшиеся деньги у пользователя
            \item \texttt{«Недостаточно денег» (str)} -  если денег не хватает
        \end{enumerate}
    
    
    \subsection*{Примеры вызова:}
        \begin{minted}{python}
        >>> bought_products(100, 10, 5)
        50
        \end{minted}

        \begin{minted}{python}
        >>> bought_products(100, 100, 5) 
        Недостаточно денег
        \end{minted}
        

% --------------------------------------------------------

    \hrulefill

    \subsection{Функция is\_in\_list}
    
    Делает проверку на наличие элемента в заранее заготовленном списке, состоящим из дней недель.
    
    \begin{minted}{python}
    def is_in_list(day):
        return day.lower() in ["сб", "вс", "суббота", "воскресенье"]
    \end{minted}
    
    \subsection*{Параметры}
        \begin{enumerate}
            \item \texttt{day (str)} - Элемент, наличие которого мы хотим проверить (День недели)
            \item \texttt{arr = ["сб", "вс", "суббота", "воскресенье"] (list): } - Массив дней недели
        \end{enumerate}
        
        \subsection*{Возвращает}
        \begin{enumerate}
            \item \texttt{True (bool)} -  Элемент есть в массиве
            \item \texttt{False (bool)} - Элемента в массиве нет        \end{enumerate}
    
    
    \subsection*{Примеры вызова:}
        \begin{minted}{python}
        >>> is_in_list("Суббота")
        True
        \end{minted}

        \begin{minted}{python}
        >>> is_in_list("С")
        False
        \end{minted}

% --------------------------------------------------------

    \hrulefill
    
    \subsection{Функция is\_even}
    
    Функция делает проверку на то, чётное ли число, и выводит результат
    
    \begin{minted}{python}
    def is_even(parametr):
        return parametr % 2 == 0
    \end{minted}
    
    \subsection*{Параметры}
        \begin{itemize}
            \item \texttt{parametr (int)} - Элемент, который мы хотим проверить
        \end{itemize}
        
        \subsection*{Возвращает}
        \begin{enumerate}
            \item \texttt{True (bool)} -  Число кратно 2
            \item \texttt{False (bool)} - Число не кратно 2
        \end{enumerate}
    
    
    \subsection*{Примеры вызова:}
        \begin{minted}{python}
    >>> is_even(2)
    True
        \end{minted}

        \begin{minted}{python}
    >>> is_even(2)
    False
        \end{minted}

% --------------------------------------------------------

    \hrulefill

    \subsection{Функция multiples\_of\_3}
    
    Функция возвращает список чисел, кратных трём
    
    \begin{minted}{python}
    def multiples_of_3(prm1, prm2, step = 1):
        res = []
        for i in range(prm1, prm2, step):
            if i % 3 == 0:
                res.append(i)
        return res
    \end{minted}
    
    \subsection*{Параметры}
        \begin{enumerate}
            \item \texttt{prm1 (int) } - Начало диапозона перебора индексов 
            \item \texttt{prm2 (int)} - Конец диапозона перебора индексов
            \item \texttt{step (int)} -  Шаг перебора индексов
                \begin{itemize}
                    \item Если при вызове ф-ции шаг \textit{(step)} \textbf{не обозначаем}: step = 1
                    \item Если при вызове ф-ции шаг  \textit{(step)} \textbf{обозначаем}: step = <введённое число>
                \end{itemize}
        \end{enumerate}
        
        \subsection*{Возвращает}
        \begin{itemize}
            \item \texttt{list} - Список чисел, которые делятся на 3  
        \end{itemize}
    
    
    \subsection*{Примеры вызова:}
        \begin{minted}{python}
    >>> multiples_of_3(0, 13) 
    [0, 3, 6, 9]
        \end{minted}

        \begin{minted}{python}
    >>> multiples_of_3(0, 13, 2) 
    [0, 6, 12]
        \end{minted}

        \begin{minted}{python}
    >>> multiples_of_3(0, 13, 2)
    [0, 6, 12]
        \end{minted}

% --------------------------------------------------------

% --------------------------------------------------------

    \hrulefill

    \subsection{Функция sentiment\_score}
    
    Функция возвращает коэффицент (оценку) того, на сколько сильно текст был эмоционально окрашен

    
    \begin{minted}{python}
    def sentiment_score(txt, pos, neg):
        words = txt.lower().split()
        ratio = 0
        
        for word in words:
            if word in pos:
                ratio += 1  
            elif word in neg:
                ratio -= 1 
    
            if word.endswith('!'): 
                ratio *= 1.5  
            elif word.endswith('?'):
                ratio *= 0.3 
                
        return round(ratio, 1)

    \end{minted}
    
    \subsection*{Параметры}
        \begin{enumerate}
            \item \texttt{txt (str)} - Строка, которую будет оценивать функция
            \item \texttt{pos (list)} - Список положительных слов (прилагательных)
            \item \texttt{neg (list)} -  Список негативных слов (прилагательных)
        \end{enumerate}
        
        \subsection*{Возвращает}
        \begin{itemize}
            \item \texttt{float} - Коэффицент эмоциональной окраски текста, округлённый до десятых  
        \end{itemize}
    
    
    \subsection*{Примеры вызова:}
        \begin{minted}{python}
    >>> sentiment_score("I love Python!", ["love", "like", "enjoy", "adore", "appreciate"], ["hate", "dislike", "detest", "antipathy"]) 
    1.5
        \end{minted}

        \begin{minted}{python}
    >>> sentiment_score("love like enjoy", ["love", "like", "enjoy"], ["hate"]) 
    3.0
        \end{minted}

        \begin{minted}{python}
    >>> sentiment_score("Do I love! Python", ["love"], ["hate"]) 
    0.0
        \end{minted}

% --------------------------------------------------------

% --------------------------------------------------------

    \hrulefill

    \subsection{Функция sentiment\_score}
    
    Функция делает перевод числа из десятичной системы счистления в двоичную и выводит результат. В текущей версии переменная count не выводится 

    
    \begin{minted}{python}
    def transform_into_binary(num, count = 0):  
        res = [] 
        
        while num > 0: 
            res.append(num % 2) 
            num //= 2 
            count += 1 
            
        if count > 10: 
            res.reverse() 
            
        return ''.join(map(str, res)) if res else "0" 
    \end{minted}
    
    \subsection*{Параметры}
        \begin{enumerate}
            \item \texttt{num (int)} - Число в десятичной системы счистления 
            \item \texttt{count (int)} - Список положительных слов (прилагательных)
            \item \texttt{neg (list)} -  Количество итераций
        \end{enumerate}
        
        \subsection*{Возвращает}
        \begin{itemize}
            \item \texttt{str} - Строка из символов входящего числа в двоичной системе счистления
        \end{itemize}
    
    
    \subsection*{Примеры вызова:}
        \begin{minted}{python}
    >>>transform_into_binary(10) 
    0101
        \end{minted}

        \begin{minted}{python}
    >>> transform_into_binary(12) 
    ('1100', 4)

        \end{minted}

        \begin{minted}{python}
    >>> transform_into_binary(12) 
    ('1100', 4)
        \end{minted}

        \begin{minted}{python}
    >>> transform_into_binary(12, 23) 
    ('1100', 27)
        \end{minted}

% --------------------------------------------------------

\hrulefill

\subsection{Граф истории коммитов}

    На рисунке показана история коммитов в ветке \texttt{doc\_561160}.
    
    \begin{figure}[h]
        \centering
        \includegraphics[width=1.3\textwidth]{log.png}
        \caption{История коммитов в ветке doc\_561160}
    \end{figure}

% --------------------------------------------------------

\hrulefill

\section{Проблемы и трудности}
    \subsection*{Информация}
        Поскольку с \textbf{LaTeX} и \textbf{Markdown} до написания лабораторной работы я не был знаком, поэтому всю информацию нужно было искать самостоятельно и тщательно делать выборку из большого количества источников. Также, я принял решение не только читать статьи на сайтах с описанием команд и документациями, но и обращаться к \textbf{LLM Hangzhou DeepSeek Artificial Intelligence Basic Technology Research Co., Ltd}, чтобы облегчить поиск нужных команд или для понимания нюансов работы с \textbf{LaTeX} и \textbf{Markdown}. Более того, на сайте https://ru.overleaf.com/ редактор предлагает вставить определённые теги, что позволет "поиграться" и попробовать в действии их посмотреть из-за встроенного на сайте компилятора для \textbf{LaTeX}, что значительно упрощало поиск подходящих тегов.
    

% --------------------------------------------------------

\hrulefill

\section{Выводы}
    
    \large{При написании данной лабораторной работы я сталкивался со множеством трудностей, начиная с понимания того, что делал имел в виду старый разраб, когда создавал эти функции, заканчивая правильностью использования тех или иных тегов при работе с \textbf{LaTeX} и \textbf{Markdown}. Благодарая этой работе я познакомился с языками разметки \textbf{LaTeX} и \textbf{Markdown} и выучил меня базовым теги, которые пригодятся в документации реального проекта.}
    

% --------------------------------------------------------

\end{document}